% `tex.web` §722 `fetch` and §723, its failure:
%
%     begin cur_c:=character(a); cur_f:=fam_fnt(fam(a)+cur_size);
%     if cur_f=null_font then
%       @<Complain about an undefined family and set |cur_i| null@>
%
%     begin print_err(""); print_size(cur_size); print_char(" ");
%     print_int(fam(a)); print(" is undefined (character ");
%     print_ASCII(qo(cur_c)); print_char(")");
%     help4("Somewhere in the math formula just ended, you used the")
%     ("stated character from an undefined font family. For example,")
%     ("plain TeX doesn't allow \it or \sl in subscripts. Proceed,")
%     ("and I'll try to forget that I needed that character.");
%     error; cur_i:=null_character; math_type(a):=empty;
%     end
%
% The character is *dropped* — `math_type(a):=empty` — not moved elsewhere.
% This engine fell back to family 0, which set a glyph TeX does not set, out
% of a font the formula never named, and said nothing about it.
%
% §699's `print_size` names the size, not the family's text font, so the same
% character in a superscript reports `\scriptfont`.  `mlist_to_hlist` has no
% error stream of its own, so the complaints are collected there and reported
% by the engine when the formula ends — which is where §723 sits in the log,
% after everything the formula traced.
% The test is `cur_f=null_font`, not "no font at all": a family explicitly set
% to `\nullfont` is as undefined as one never set.  And §1233 reads the
% right-hand side with `scan_font_ident`, which takes any `<font>` — a
% `\font`-defined identifier, the primitive `\font` for the current one, or
% another family.  Matching only the first of those left `\textfont1=\font`
% quietly unset, which stayed invisible until §723 began reporting.
\catcode`\{=1 \catcode`\}=2 \catcode`\$=3 \catcode`\^=7
\font\r=cmr10 \r \hsize100pt \parindent0pt
\textfont2=\r \scriptfont2=\r \scriptscriptfont2=\r
\textfont3=\r \scriptfont3=\r \scriptscriptfont3=\r
\fontdimen22\r=1pt
$\fam0 A^{\fam0 B}$\par
% §722 `fetch` is called from more than the ordinary conversion.  §738's
% `make_math_accent` opens with `fetch(accent_chr(q))`, and when it fails the
% accent block is *skipped* while the nucleus is still converted:
%
%     begin fetch(accent_chr(q));
%     if char_exists(cur_i) then
%       begin i:=cur_i; c:=cur_c; f:=cur_f; … end;
%     end;
%
% And §749's `make_op` opens with `fetch(nucleus(q))` — before the limit boxes
% are cleaned, so an operator out of an undefined family is reported ahead of
% its own superscript.  §723's `math_type(a):=empty` is what stops either from
% being reported twice.
$\mathaccent"7161 {A}$\par
$\mathop{\fam0 A}\limits^{\fam0 B}$\par
$\radical"161 {\fam0 A}$\par
\textfont1=\font \showthe\textfont1
\textfont0=\nullfont \scriptfont0=\nullfont \scriptscriptfont0=\nullfont
$\fam0 C$\par
\end
